% Supported v0.5 statement-of-purpose example.
\documentclass[type=purpose]{careerdossier-statement}
% examples/profiles/profile-academic.tex
%
% Shared profile metadata used by the academic-CV example and the corresponding
% industry-resume smoke fixture. Keeping it independent of document layout
% proves that both document families consume the same \CDossierSetup interface.

\CDossierSetup{
  name     = {Ada Lovelace},
  headline = {Researcher in Analytical Computing},
  email    = {ada.lovelace@example.com},
  location = {Toronto, Ontario, Canada},
  website  = {example.com/ada-lovelace},
  %% The bare Google Scholar identifier, expanded to
  %% scholar.google.com/citations?user=ada-example for both the displayed text
  %% and the link. Pasting the whole address works too and is left as written.
  %% See docs/API.md, "Accepted forms for the web-profile keys".
  scholar  = {ada-example},
  orcid    = {0000-0002-1825-0097},
  affiliation = {Example University},
}

\CDossierStatementSetup{
  subtitle={Preparation and goals in computational science},
  application-context={Application to the MSc in Computational Science},
  application-id={12345678}
}
\begin{document}
\MakeCDossierStatementHeader

\section*{Purpose and Motivation}

I am applying to the MSc in Computational Science at Example University to
develop the mathematical and research foundations needed to build reliable
scientific software. My experience has shown me that computational work is most
useful when a result can be explained, tested, and reproduced by someone beyond
the original project. I want advanced training that connects numerical methods,
software design, and the interpretation of evidence rather than treating them
as separate technical skills.

This goal emerged while I was working on small data-analysis projects in which
apparently minor choices produced materially different conclusions. A change in
missing-data handling, a default tolerance, or an undocumented software version
could alter an output without making the source of the change visible. I became
interested not only in obtaining an answer, but in understanding which parts of
the process made that answer defensible.

\section*{Academic Preparation}

My undergraduate work in mathematics gave me a foundation in analysis, linear
algebra, probability, and differential equations. Courses in numerical analysis
were especially influential because they joined theoretical questions about
error and convergence with practical questions about representation and finite
resources. I learned to evaluate an algorithm through both proof and experiment,
and to distinguish a method's mathematical properties from the behaviour of a
particular implementation.

Alongside this coursework, I developed proficiency in programming through
projects that required data preparation, simulation, visualization, and written
interpretation. I began using version control and automated tests after a result
could not be reproduced from an earlier project directory. Rebuilding that work
made the value of explicit environments and small verification cases immediate.
These practices now shape how I approach every computational assignment.

My capstone project examined adaptive approximation for a model with spatially
varying parameters. I was responsible for implementing the solver, constructing
reference cases, and explaining the relationship between estimated and observed
error. The initial implementation produced plausible figures but failed on a
boundary case. Tracing that failure led me to redesign the tests and to document
the assumptions connecting the model and code. The final project was stronger
because the unexpected result became part of the analysis rather than something
to hide.

\section*{Relevant Experience and Growth}

As a peer tutor, I have supported students in introductory programming and
calculus. Explaining the same idea through diagrams, algebra, and executable
examples helped me recognize the difference between performing a procedure and
understanding why it works. Tutoring also taught me to listen for the reasoning
behind an error before offering a correction. I expect this experience to
support both collaborative coursework and future teaching-assistant duties.

I have also contributed documentation and tests to a small open-source
scientific project. Working with code written by several people required me to
read design decisions from incomplete records, ask focused questions, and make
changes that did not disrupt existing users. This experience strengthened my
interest in research software as a scholarly object whose quality affects the
credibility and reach of scientific work.

\section*{Programme Fit and Goals}

Example University's programme is a strong fit because it combines advanced
numerical study with interdisciplinary research and an explicit emphasis on
reproducible computation. I am particularly interested in courses on numerical
linear algebra, uncertainty quantification, and research-software engineering.
The opportunity to work with faculty and students across applied mathematics
and domain sciences would let me test methods against questions that cannot be
understood from a single disciplinary perspective.

During the degree, I hope to investigate verification strategies for pipelines
that combine simulation and data-driven components. My near-term goal is to
develop a thesis project with a clear mathematical question, an open reference
implementation, and an evaluation design that makes limitations visible. After
the MSc, I intend to continue in research or research-software work where I can
help interdisciplinary teams produce computational evidence that others can
inspect and reuse.

I would bring careful preparation, experience learning from failure, and a
commitment to shared technical practice. The programme would give me the depth
and mentorship to turn those habits into a coherent research direction and to
contribute more responsibly to computational science.

\end{document}
